% Regression: hard09_pull — exercises T09 — RMP sptintertwin-gh: pulling-through / intertwiner equation. An MPO…
% Formula: geometry-only; no tensor identity.
% T09 — RMP sptintertwin-gh: pulling-through / intertwiner equation.
% An MPO (red operator wire) crosses the virtual wire inside a tall box
% X_gh, and single-site operators g, h sit ON the physical legs.
\documentclass[varwidth=19cm, border=6pt]{standalone}
\usepackage{amsmath, amssymb}
\usepackage{tenkz}
\begin{document}

(a) grid attempt — operator row above a ket row: the red wire lives on its
own row, so it can never CROSS the black virtual wire inside the tall box, and g/h
cannot sit on the vertical paired legs:
\[
  \begin{tenkz}[rows={op:operator, ket}, tensor style=box]
    \tnX{gh} & \tn[box]{X_{gh}} & \tn[box]{g} & \tn[box]{h}\\
             & \tn{A} & \tn{A} & \tn{A}
  \end{tenkz}
\]

(b) free-tier attempt — raw coordinates; the crossing and the on-leg boxes
are hand-placed; the red wire hue is not reachable from the free tier:
\[
  \begin{tenkzfree}
    \tnput[ring]{gh}{(-1.1,0.35)}{gh}
    \tnput[box]{X}{(0,0)}{X_{gh}}
    \tnput[dot]{a1}{(1.1,-0.35)}{} \tnput[dot]{a2}{(2.2,-0.35)}{}
    % black virtual wire through X
    \tnjoin{-1.9,-0.35}{-0.6,-0.35}
    \tnjoin{-0.6,-0.35}{X.west}
    \tnjoin{X.east}{a1.west} \tnjoin{a1.east}{a2.west}
    \tnjoin{a2.east}{2.9,-0.35}
    % red MPO wire entering X high and leaving low (the crossing)
    \tnjoin{-1.9,0.35}{gh.west} \tnjoin{gh.east}{X.140}
    \tnjoin{X.south}{0,-0.9}
    % physical legs with on-leg boxes g, h
    \tnput[box]{g}{(1.1,0.35)}{g} \tnput[box]{h}{(2.2,0.35)}{h}
    \tnjoin{a1.north}{g.south} \tnjoin{g.north}{1.1,0.9}
    \tnjoin{a2.north}{h.south} \tnjoin{h.north}{2.2,0.9}
  \end{tenkzfree}
\]
\end{document}
